\section{深度学习时代：计算机视觉的爆发}

\subsection{视觉识别任务的全面突破}

2012 年之后，深度学习迅速在计算机视觉的各个任务上取得了突破性进展：

\begin{figure}[H]
\centering
\includegraphics[width=0.95\textwidth]{slides-images/slide-050.jpg}
\caption{深度学习在视觉任务中的应用：左侧为目标检测（Object Detection，Ren et al., 2015），右侧为图像语义分割（Image Segmentation，Fabaret et al., 2012）\protect\footnotemark}
\end{figure}
\footnotetext{Slides 第51页（Part 1）。}

\begin{itemize}
    \item \textbf{物体检测}（Object Detection）：不仅识别``图中有什么''，还定位``在哪里''
    \item \textbf{图像分割}（Image Segmentation）：对每个像素进行分类
    \item \textbf{视频分类}（Video Classification）：理解视频中的动作和活动
    \item \textbf{图像描述}（Image Captioning）：自动为图像生成文字描述（Andrej Karpathy 在 Fei-Fei Li 实验室的博士论文工作）
    \item \textbf{关系理解}（Relationship Understanding）：理解物体之间的空间和语义关系
    \item \textbf{风格迁移}（Style Transfer）：将一种视觉风格应用到另一张图像上
\end{itemize}

\begin{figure}[H]
\centering
\includegraphics[width=0.95\textwidth]{slides-images/slide-055.jpg}
\caption{关系理解：检测图像中物体之间的空间关系（spatial）、比较关系（comparative）和动作关系（verb），构建场景图（Scene Graph）\protect\footnotemark}
\end{figure}
\footnotetext{Slides 第56页（Part 1）。来源：Krishna, Lu, Bernstein, Fei-Fei, ECCV 2016。}

\subsection{三大驱动力的汇聚}

Fei-Fei Li 总结了推动当前 AI 爆发的三大汇聚力量：

\begin{figure}[H]
\centering
\includegraphics[width=0.95\textwidth]{slides-images/slide-060.jpg}
\caption{AI 爆发的三大驱动力：计算力（Computation）、算法（Algorithms）、数据（Data）\protect\footnotemark}
\end{figure}
\footnotetext{Slides 第61页（Part 1）。}

\begin{importantbox}{AI 爆发的三大驱动力}
\begin{enumerate}
    \item \textbf{计算力}（Computation）：GPU 性能的指数级增长（2020年后加速尤为明显），NVIDIA GPU 的 GFLOPS/dollar 大幅提升
    \item \textbf{算法}（Algorithms）：从 CNN 到 Transformer，不断涌现的网络架构创新
    \item \textbf{数据}（Data）：互联网和数字相机提供了海量训练数据
\end{enumerate}
这三者的同时就绪，将 AI 从``寒冬''推向了 Fei-Fei Li 所说的``AI Global Warming''——一个目前看不到减速迹象的爆发时代。
\end{importantbox}

\subsection{生成式 AI 时代}

课程也展示了生成式 AI 的最新进展：

\begin{itemize}
    \item \textbf{人脸生成}：使用 GAN 等技术生成逼真的人脸
    \item \textbf{文生图}（Text-to-Image）：DALL-E、Midjourney 等模型可以根据文字描述生成图像
    \item \textbf{视频理解与生成}：多模态模型结合视觉、语音等多种信息源
\end{itemize}

\subsection{应用领域}

深度学习驱动的计算机视觉已经在众多领域产生了深远影响：

\begin{itemize}
    \item \textbf{医学影像}：放射学、病理学等视觉密集型医学领域
    \item \textbf{科学发现}：如首张黑洞照片的计算成像
    \item \textbf{环境与可持续发展}：遥感监测、生态保护
    \item \textbf{自动驾驶}：道路场景理解、行人检测
    \item \textbf{机器人学}：具身智能、操作和导航
\end{itemize}

\subsection{本章小结}

深度学习在 2012 年后彻底改变了计算机视觉的面貌。从目标检测到图像分割，从图像描述到生成式 AI，深度神经网络在几乎所有视觉任务上都取得了突破。计算力、算法和数据三大力量的汇聚，将 AI 推入了一个前所未有的爆发期。

%% ========== 伦理与社会影响 ==========

